% from ega3-6-fr.tex
\oldpage[III]{170}
par la méthode de (6.6.2), à l'aide de l'image réciproque d'un recouvrement
affine fixe $\mathfrak U^{(i)}$ de $X^{(i)}$ ($1\leq i\leq m$)
(indépendant de $S'$ et des $\mathscr P_\bullet^{(i)}$); les termes
$L_{jk}^{(i)}$ sont nuls pour $j<-N_i$ (ne dépendant que de
$\mathfrak U^{(i)}$); si l'un des $q_i$ n'est pas nul, le complexe simple dont
l'homologie de degré $p$ est (6.9.7.1) a ses termes nuls pour tous les degrés
\[
  <r+\sum_{j\ne i}d_j-\sum_{i=1}^m N_i,
\]
donc (6.9.7.1) est nul pour $p<r-N$, en désignant par $N$ le plus grand des
nombres
\[
  \sum_{i=1}^m N_i-\sum_{j\ne i}d_j.
\]
On en conclut que l'on a ${}^{(f')}\mathscr E_{pq}^2=0$ pour $q\ne0$ et
$p<r-N$; comme d'autre part ${}^{(f')}\mathscr E_{pq}^2=0$ pour $q<0$, on
voit que le edge-homomorphisme (6.9.6.4) est bijectif pour $n<r-N$
(M, XV, 5.6) (pour $\mathscr Q'_\bullet=\mathscr O_{S'}$). En second lieu, si
\[
  \mathfrak{Tor}_q^S
  (f_1,\ldots,f_m;
   \mathscr P_\bullet^{(1)},\ldots,\mathscr P_\bullet^{(m)})
\]
est plat sur $S$ pour $q<r$, on a ${}^{(e)}\mathscr E_{pq}^2=0$ pour $p\ne0$
et $q<r$; par ailleurs ${}^{(e)}\mathscr E_{pq}^2=0$ pour $p<0$, donc le
edge-homomorphisme (6.9.6.5) est bijectif pour $n<r$, ce qui achève la
démonstration.
\end{proof}

Le cas le plus important de (6.9.3) dans les applications est celui où $m=1$,
$\mathscr Q'_\bullet$ étant réduit à un seul terme $\mathscr F'$ de degré
$0$; nous l'énoncerons à nouveau dans ce cas en vue de références
ultérieures\footnote{Le cas traité dans (6.9.8), et en particulier les suites
spectrales (6.9.8.3) et (6.9.8.4), nous avaient été signalés en 1957 par
J.-P. Serre.} :

\begin{proposition}[6.9.8]
Soient $S$ un préschéma, $g:S'\to S$ un morphisme,
$f:X\to Y$ un $S$-morphisme séparé et quasi-compact de $S$-préschémas,
$\mathscr P_\bullet$ un complexe de $\mathscr O_X$-Modules quasi-cohérents
limité inférieurement, $\mathscr F'$ un $\mathscr O_{S'}$-Module
quasi-cohérent. Il existe deux foncteurs spectraux biréguliers en
$\mathscr P_\bullet$ et $\mathscr F'$, à valeurs dans la catégorie des
$\mathscr O_{Y_{(S')}}$-Modules quasi-cohérents, ayant même aboutissement
$\mathfrak{Tor}_\bullet^S(f,1_{S'};\mathscr P_\bullet,\mathscr F')$, et dont
les termes $E_2$ sont
\[
  {}'\mathscr E_{pq}^2
  =\mathscr{T}\!or_p^S
  (\mathscr H^{-q}(f,\mathscr P_\bullet),\mathscr F')
  \tag{6.9.8.1}\label{III.6.9.8.1.fr}
\]
et
\[
  {}''\mathscr E_{pq}^2
  =\mathscr H^{-p}
  (f',\mathscr{T}\!or_q^S(\mathscr P_\bullet,\mathscr F')),
  \tag{6.9.8.2}\label{III.6.9.8.2.fr}
\]
où $f'=f_{(S')}:X_{(S')}\to Y_{(S')}$. 
\end{proposition}

\begin{proof}
Les suites en question peuvent aussi s'obtenir, non en partant de (6.9.3),
mais des suites $(a)$ et $(b')$ de (6.7.3) pour
$X^{(1)}=X$, $Y^{(1)}=Y$, $X^{(2)}=Y^{(2)}=S'$, $f_1=f$,
$f_2=1_{S'}$. Lorsque $S=S'=Y$, $Y$ étant affine, on obtient deux suites
spectrales de termes $E_2$ égaux à
\[
  {}'\mathscr E_{pq}^2
  =\mathscr{T}\!or_p^Y
  (\mathscr H^{-q}(f,\mathscr P_\bullet),\mathscr F)
  \tag{6.9.8.3}\label{III.6.9.8.3.fr}
\]
\[
  {}''\mathscr E_{pq}^2
  =\mathscr H^{-p}
  (f,\mathscr{T}\!or_q^Y(\mathscr P_\bullet,\mathscr F))
  \tag{6.9.8.4}\label{III.6.9.8.4.fr}
\]
aboutissant (en vertu de (6.7.6)) à l'hypercohomologie
$\mathscr H^\bullet(f,\mathscr P_\bullet\otimes_Y\mathscr F)$ du foncteur
$f_*$ par rapport au complexe
$\mathscr P_\bullet\otimes_Y\mathscr F$ de $\mathscr O_X$-Modules, pour tout
$\mathscr O_Y$-Module quasi-cohérent et $Y$-plat $\mathscr F$ (ou pour tout
$\mathscr O_Y$-Module quasi-cohérent $\mathscr F$ lorsque
$\mathscr P_\bullet$ est formé de $\mathscr O_X$-Modules $Y$-plats), qui sont
distinctes de celles de (6.2.1).
\end{proof}

\begin{corollary}[6.9.9]
Sous les conditions de (6.9.8), supposons que le complexe
$\mathscr P_\bullet$ soit limité inférieurement, formé de Modules plats sur
$S$, et que les $\mathscr O_Y$-Modules
$\mathscr H^n(f,\mathscr P_\bullet)$ soient plats sur $S$.

